\clearpage
\VerdierPhysicalPageAnchor{0149}%
\begingroup
\small
\setlength{\parskip}{0pt}%
\setlength{\emergencystretch}{1.5em}%
\setlength{\abovedisplayskip}{3pt plus 1pt minus 1pt}%
\setlength{\belowdisplayskip}{3pt plus 1pt minus 1pt}%
\sloppy

\VerdierSemanticAnchor{ch2-2-3-5-injectivity-p149-continuation}
\noindent かつ $uf=0$ となる。a), i) を用いると、さらに可換図式
\VerdierSemanticAnchor{ch2-2-3-5-injectivity-refinement-diagram}
\begin{center}
 \begin{tikzcd}[column sep=3.7em,row sep=2.3em]
 && {} \\
 & {} \arrow[ur,"v"] && {t'\in S_{\mathcal A}(\mathcal A\cap\mathcal B)} \\
 X \arrow[ur,"f"] && Y \arrow[ul,"s"] \arrow[uu,"{t'}"'] &
 \end{tikzcd}
\end{center}
\VerdierSemanticAnchor{ch2-2-3-5-injectivity-conclusion}
\noindent が存在して $vf=0$ となる。したがって
$m=Q_1(t')^{-1}Q_1(vf)=0$ である。
\VerdierSemanticAnchor{ch2-2-3-5-surjectivity-map}
次に準同型
\[
 \operatorname{Hom}_{\mathcal A/(\mathcal A\cap\mathcal B)}(X,Y)
 \longrightarrow\operatorname{Hom}_{\mathcal D/\mathcal B}(X,Y)
\]
が全射であることを示す。
\VerdierSemanticAnchor{ch2-2-3-5-surjectivity-roof}
任意の元 $p\in\operatorname{Hom}_{\mathcal D/\mathcal B}(X,Y)$ は
\[
 p=Q_2(t)^{-1}Q_2(g),\qquad t\in S_{\mathcal D}(\mathcal B)
  .
\]
と書ける。
\VerdierSemanticAnchor{ch2-2-3-5-surjectivity-refinement}
a), i) により、この元は常に
\[
 p=Q_2(s)^{-1}Q_2(g'),\qquad
 s\in S_{\mathcal A}(\mathcal A\cap\mathcal B),
\]
と書ける。
\VerdierSemanticAnchor{ch2-2-3-5-surjectivity-conclusion}
したがって $\operatorname{Hom}_{\mathcal A/(\mathcal A\cap\mathcal B)}(X,Y)$
の元から来る。
\VerdierSemanticAnchor{ch2-2-3-5-c-proof-i-implies-ii}
残る c) の同値性を示す。まず i) $\Rightarrow$ ii) を示す。
\VerdierSemanticAnchor{ch2-2-3-5-iterated-quotient}
(2.3.1) により $\mathcal D/\mathcal B$ は標準的に
$(\mathcal D/(\mathcal A\cap\mathcal B))
/(\mathcal B/(\mathcal A\cap\mathcal B))$ と同型であり、標準関手
\VerdierSemanticAnchor{ch2-2-3-5-q-functor}
\[
 Q\colon\mathcal D/(\mathcal A\cap\mathcal B)\longrightarrow\mathcal D/\mathcal B
\]
は商関手である。
\VerdierSemanticAnchor{ch2-2-3-5-q-free-reduction}
よって命題 (2.3.3, a)) により、i) $\Rightarrow$ ii) を示すには、
$\mathcal A$ の任意の対象 $X$ と
$\mathcal B$ の任意の対象 $B$ に対して
$\operatorname{Hom}_{\mathcal D/(\mathcal A\cap\mathcal B)}(B,X)=0$ を
示せばよい。
\VerdierSemanticAnchor{ch2-2-3-5-m-representation}
$\operatorname{Hom}_{\mathcal D/(\mathcal A\cap\mathcal B)}(B,X)$ の任意の元
$m$ は $Q_3(f)Q_3(s)^{-1}$ と書ける。ここで
\VerdierSemanticAnchor{ch2-2-3-5-q3-functor}
\[
 Q_3\colon\mathcal D\longrightarrow\mathcal D/(\mathcal A\cap\mathcal B)
\]
は標準関手であり、$s$ は $S_{\mathcal D}(\mathcal A\cap\mathcal B)$ に属し、後者は
$\mathcal A\cap\mathcal B$ に対応する $\mathcal D$ の乗法的系である
(2.2.2)。よって $\mathcal D$ に次の図式がある：
\VerdierSemanticAnchor{ch2-2-3-5-w-diagram}
\begin{center}
 \begin{tikzcd}[column sep=4em,row sep=2.4em]
 & W \arrow[dl,"s"'] \arrow[dr,"f"] & \\
 B && X
 \end{tikzcd}\quad ,
\end{center}
\endgroup
